% ============================================================
\chapter{Variational Methods}
\label{ch:variational}
% ============================================================

\section{Weak formulation of elliptic problems}

\begin{definition}[Associated bilinear form]
\index{bilinear form}
Consider the Dirichlet problem for a second-order elliptic PDE:
\begin{equation}\label{eq:dirichlet}
    -\sum_{i,j=1}^n \partial_j(a_{ij}(x) \partial_i u) + c(x) u = f(x)
    \quad \text{in } \Omega, \qquad u = 0 \quad \text{on } \partial\Omega.
\end{equation}
The \emph{weak formulation} consists of finding $u \in H_0^1(\Omega)$ such that:
\[
    a(u, v) = \ell(v) \quad \forall\, v \in H_0^1(\Omega),
\]
where the bilinear form $a$ and the linear form $\ell$ are defined by:
\begin{align*}
    a(u, v) &= \int_\Omega \left( \sum_{i,j} a_{ij} \partial_i u \, \partial_j v
    + c\, u\, v \right) dx, \\
    \ell(v) &= \int_\Omega f\, v\, dx.
\end{align*}
\end{definition}

\begin{intuition}
The weak formulation is obtained by multiplying the PDE by a test function
$v \in H_0^1(\Omega)$ and integrating by parts. The advantage is twofold: the
order of differentiation required on $u$ is reduced (from $2$ to $1$), and one
obtains a functional framework where theorems of functional analysis apply
directly.
\end{intuition}

\begin{definition}[Coercivity and continuity]
\index{coercivity}
Let $V$ be a Hilbert space. A bilinear form $a : V \times V \to \R$ is:
\begin{itemize}
    \item \textbf{continuous} if there exists $M > 0$ such that $\abs{a(u,v)} \leq M\norm{u}_V \norm{v}_V$;
    \item \textbf{coercive} if there exists $\alpha > 0$ such that $a(u,u) \geq \alpha \norm{u}_V^2$.
\end{itemize}
\end{definition}

\section{The Lax-Milgram theorem}

\begin{theorem}[Lax-Milgram]
\label{thm:lax_milgram}
\index{Lax-Milgram theorem}
Let $V$ be a Hilbert space, $a : V \times V \to \R$ a continuous and coercive
bilinear form, and $\ell : V \to \R$ a continuous linear form. Then there exists a
unique $u \in V$ such that:
\[
    a(u, v) = \ell(v) \quad \forall\, v \in V.
\]
Moreover, $\norm{u}_V \leq \frac{1}{\alpha}\norm{\ell}_{V'}$.
\end{theorem}

\begin{proof}
\textbf{Step 1: Riesz representation.}
By the Riesz representation theorem, for each fixed $u \in V$, there exists a
unique $Au \in V$ such that $a(u, v) = \inner{Au}{v}_V$ for all $v$. The map
$A : V \to V$ is linear and continuous with $\norm{A} \leq M$.

Similarly, there exists $F \in V$ such that $\ell(v) = \inner{F}{v}_V$.

The problem reduces to: find $u$ such that $Au = F$.

\textbf{Step 2: $A$ is invertible.}
Coercivity gives $\inner{Au}{u} = a(u,u) \geq \alpha\norm{u}^2$, so $A$ is
injective and $\norm{Au} \geq \alpha\norm{u}$. The image $\operatorname{Im}(A)$
is closed (since $A$ is bounded below). If $w \perp \operatorname{Im}(A)$, then
$\inner{Aw}{w} = 0$, hence $w = 0$ by coercivity. Thus
$\operatorname{Im}(A) = V$ and $A$ is bijective. The unique solution is
$u = A^{-1}F$ with $\norm{u} \leq \frac{1}{\alpha}\norm{F} = \frac{1}{\alpha}\norm{\ell}_{V'}$.
\end{proof}

\begin{warning}[title=\textbf{Symmetry not required}]
The Lax-Milgram theorem does not require $a$ to be symmetric. This is an advantage
over the Riesz theorem, which would require symmetry.
\end{warning}

\section{Application to the Dirichlet problem}

\begin{theorem}[Existence and uniqueness for the Dirichlet problem]
\label{thm:dirichlet}
Let $\Omega \subset \R^n$ be a bounded open set, $f \in L^2(\Omega)$, and suppose
the coefficients $a_{ij} \in L^\infty(\Omega)$ satisfy the uniform ellipticity
condition:
\[
    \sum_{i,j} a_{ij}(x) \xi_i \xi_j \geq \lambda \abs{\xi}^2
    \quad \forall\, \xi \in \R^n,\; \text{a.e.\ } x \in \Omega,
\]
and $c(x) \geq 0$ a.e. Then problem \eqref{eq:dirichlet} has a unique weak
solution $u \in H_0^1(\Omega)$.
\end{theorem}

\begin{proof}
We verify the hypotheses of Lax-Milgram with $V = H_0^1(\Omega)$.

\textbf{Continuity of $a$:}
$\abs{a(u,v)} \leq \norm{a_{ij}}_{L^\infty} \norm{\nabla u}_{L^2} \norm{\nabla v}_{L^2}
+ \norm{c}_{L^\infty} \norm{u}_{L^2} \norm{v}_{L^2} \leq M \norm{u}_{H^1} \norm{v}_{H^1}$.

\textbf{Coercivity:}
$a(u,u) \geq \lambda \norm{\nabla u}_{L^2}^2 \geq \frac{\lambda}{1 + C_P^2} \norm{u}_{H^1}^2$
by the Poincar\'e inequality (where $C_P$ is the Poincar\'e constant).

\textbf{Continuity of $\ell$:}
$\abs{\ell(v)} = \abs{\int f v} \leq \norm{f}_{L^2}\norm{v}_{L^2}
\leq \norm{f}_{L^2}\norm{v}_{H^1}$.

By Lax-Milgram, there exists a unique solution $u \in H_0^1(\Omega)$.
\end{proof}

\begin{example}[Laplacian with source]
For $-\Delta u = f$ in $\Omega$, $u = 0$ on $\partial\Omega$:
$a(u,v) = \int_\Omega \nabla u \cdot \nabla v\, dx$ and $\ell(v) = \int_\Omega fv\, dx$.
The hypotheses are satisfied with $\lambda = 1$ and $c = 0$.
\end{example}

\section{The Ritz-Galerkin method}

\begin{definition}[Galerkin approximation]
\index{Galerkin method}
Let $V_h \subset V$ be a finite-dimensional subspace. The \emph{Galerkin
approximation} of $u$ is the unique $u_h \in V_h$ such that:
\[
    a(u_h, v_h) = \ell(v_h) \quad \forall\, v_h \in V_h.
\]
\end{definition}

\begin{theorem}[C\'ea's lemma]
\label{thm:cea}
\index{C\'ea's lemma}
If $a$ is continuous (constant $M$) and coercive (constant $\alpha$), then:
\[
    \norm{u - u_h}_V \leq \frac{M}{\alpha} \inf_{v_h \in V_h} \norm{u - v_h}_V.
\]
\end{theorem}

\begin{proof}
By coercivity and the Galerkin equation (orthogonality of the error):
\begin{align*}
    \alpha\norm{u - u_h}^2 &\leq a(u - u_h, u - u_h) \\
    &= a(u - u_h, u - v_h) + a(u - u_h, v_h - u_h) \\
    &= a(u - u_h, u - v_h) \quad \text{(since $a(u - u_h, w_h) = 0$ for $w_h \in V_h$)} \\
    &\leq M\norm{u - u_h}\norm{u - v_h}.
\end{align*}
Dividing by $\norm{u - u_h}$ and taking the infimum over $v_h$ gives the result.
\end{proof}

\begin{remark}
C\'ea's lemma shows that the Galerkin approximation is quasi-optimal: the error
is at most $M/\alpha$ times the best possible approximation in $V_h$.
\end{remark}

\section{Variational principle (symmetric case)}

\begin{theorem}[Variational equivalence]
\label{thm:variational_equivalence}
If $a$ is symmetric, continuous and coercive, then $u$ solves $a(u,v) = \ell(v)$
for all $v \in V$ if and only if $u$ minimizes the energy functional:
\[
    J(v) = \frac{1}{2} a(v, v) - \ell(v).
\]
\end{theorem}

\begin{proof}
For any $v \in V$:
$J(v) = J(u) + \frac{1}{2}a(v - u, v - u) + a(u, v - u) - \ell(v - u)$.
If $u$ is a solution, the last two terms cancel, giving
$J(v) = J(u) + \frac{1}{2}a(v-u, v-u) \geq J(u)$ by coercivity.
Conversely, if $u$ minimizes $J$, then $\frac{d}{dt}J(u + tv)|_{t=0} = 0$
gives $a(u,v) = \ell(v)$.
\end{proof}

\begin{example}[Dirichlet principle]
\index{Dirichlet principle}
For the problem $-\Delta u = f$ with $u \in H_0^1(\Omega)$, the energy functional is:
\[
    J(v) = \frac{1}{2}\int_\Omega \abs{\nabla v}^2\, dx - \int_\Omega f v\, dx.
\]
Minimizing $J$ is equivalent to solving $-\Delta u = f$.
\end{example}

\section{Regularity of weak solutions}

\begin{theorem}[$H^2$ regularity]
\label{thm:H2_regularity}
\index{elliptic regularity}
Let $\Omega$ be a bounded open set with $C^2$ boundary and $u \in H_0^1(\Omega)$
the weak solution of $-\Delta u = f$ with $f \in L^2(\Omega)$. Then
$u \in H^2(\Omega)$ and:
\[
    \norm{u}_{H^2(\Omega)} \leq C\norm{f}_{L^2(\Omega)}.
\]
\end{theorem}

\begin{theorem}[$C^\infty$ regularity]
If $f \in C^\infty(\overline{\Omega})$ and $\partial\Omega$ is $C^\infty$, then
the weak solution $u \in H_0^1(\Omega)$ of $-\Delta u = f$ satisfies
$u \in C^\infty(\overline{\Omega})$.
\end{theorem}

\begin{remark}
Regularity depends on the smoothness of the boundary and the right-hand side. If
$\partial\Omega$ is not smooth enough (e.g., a domain with corners), $H^2$
regularity may fail.
\end{remark}

\begin{keyformulas}[title=\textbf{Summary of the variational method}]
\begin{enumerate}
    \item Write the weak formulation: $a(u,v) = \ell(v)$, $\forall v \in V$.
    \item Verify continuity of $a$ and $\ell$.
    \item Verify coercivity of $a$ (often via Poincar\'e).
    \item Apply Lax-Milgram $\Rightarrow$ existence and uniqueness.
    \item Study regularity: $u \in H^2$? $u \in C^\infty$?
\end{enumerate}
\end{keyformulas}

\section{Exercises}

\begin{exercise}
Write the weak formulation of the problem:
$-\operatorname{div}(A(x)\nabla u) + u = f$ in $\Omega$, $u = 0$ on $\partial\Omega$,
where $A(x)$ is a uniformly positive definite symmetric matrix. Verify the
hypotheses of Lax-Milgram.
\end{exercise}

\begin{exercise}
Show that for the Neumann problem:
$-\Delta u = f$ in $\Omega$, $\frac{\partial u}{\partial \nu} = 0$ on
$\partial\Omega$, the compatibility condition $\int_\Omega f\, dx = 0$ is
necessary. Formulate the weak problem in $H^1(\Omega)$ and apply Lax-Milgram
in the quotient $H^1/\R$.
\end{exercise}

\begin{exercise}
Consider the Robin problem: $-\Delta u + u = f$ in $\Omega$,
$\frac{\partial u}{\partial \nu} + \alpha u = 0$ on $\partial\Omega$ with
$\alpha > 0$. Write the weak formulation in $H^1(\Omega)$ and prove existence
and uniqueness.
\end{exercise}

\begin{exercise}[C\'ea's lemma in practice]
Let $\Omega = (0,1)$, $-u'' = f$ with $u(0) = u(1) = 0$. Take $V_h$ to be the
space of piecewise linear functions on a uniform mesh of size $h$. Show that
$\norm{u - u_h}_{H^1} \leq C h \norm{u}_{H^2}$.
\end{exercise}

\begin{exercise}
Show that if $a$ is symmetric, the Galerkin approximation $u_h$ minimizes
$J(v) = \frac{1}{2}a(v,v) - \ell(v)$ over $V_h$. Deduce that the energy error
satisfies $a(u - u_h, u - u_h) \leq a(u - v_h, u - v_h)$ for all $v_h \in V_h$.
\end{exercise}

\begin{exercise}
Let $\Omega$ be an L-shaped domain in $\R^2$. Give an example of data
$f \in L^2(\Omega)$ for which the solution of $-\Delta u = f$ is not in
$H^2(\Omega)$. Explain the role of the corner singularity.
\end{exercise}
